\chapter{Agreement, Lists and Mutual Correlated Agreement}
\label{ch:mca}

This chapter isolates the only coding-theoretic input of the soundness analysis:
\emph{mutual correlated agreement} (MCA) for the univariate powers generator
$\beta\mapsto(1,\beta,\dots,\beta^{a-1})$, which is exactly the generator of the fold
operator of \cref{def:fold}. In the unique-decoding regime MCA is proved here from the
proximity gap for curves; up to the Johnson radius it is imported. The chapter also
proves the unique block decoding lemma and the Johnson list bound for block agreement,
which are used to identify the reference polynomial in later chapters.

Throughout, $E\subseteq\K$ is a set of $n$ distinct points, $\ell\le n$,
$C:=\RS[\K,E,\ell]=\{P|_E:P\in\K[X],\deg P<\ell\}$ and $\rho:=\ell/n$.

%=====================================================================
\section{Mutual correlated agreement}
%=====================================================================

\begin{definition}[Bad challenges, MCA error]\label{def:mca}
Let $a\ge2$, $u=(u_0,\dots,u_{a-1})$ with $u_j:E\to\K$, and $\delta\in[0,1)$. A challenge
$\beta\in\K$ is \emph{bad} for $(u,\delta)$ if there exist $S\subseteq E$ with
$|S|\ge(1-\delta)n$ and $v\in C$ such that
\begin{enumerate}[label=\textup{(\roman*)}]
\item $\sum_{j<a}\beta^ju_j(x)=v(x)$ for all $x\in S$, and
\item there is no tuple $(v_0,\dots,v_{a-1})\in C^a$ with $u_j|_S=v_j|_S$ for all $j<a$.
\end{enumerate}
Let $\Bad(u,\delta)\subseteq\K$ be the set of bad challenges and
\[
\eps_{\MCA}(C,a,\delta):=\max_{u\in(\K^E)^a}\frac{|\Bad(u,\delta)|}{|\K|}.
\]
\end{definition}

\begin{remark}\label{rem:mca-monotone}
If $\delta\le\delta'$ then $\Bad(u,\delta)\subseteq\Bad(u,\delta')$, so
$\eps_{\MCA}(C,a,\cdot)$ is non-decreasing. The word ``mutual'' refers to the fact that
the set $S$ in the conclusion is the \emph{same} set on which the combination agrees with
$v$; ordinary correlated agreement only asserts the existence of some large set.
\end{remark}

The following lemma translates \cref{def:mca} into the language of folding. Fix
$0\le i<k$, $a:=a_{i+1}$, and take $E=D_{i+1}$, $C=\RS_{i+1}$, so $n=n_{i+1}$ and
$\ell=L_{i+1}$.

\begin{lemma}[MCA in folding form]\label{lem:mca-fold}
Let $w:D_i\to\K$, $u:=(u_j[w])_{j<a}$, $\delta\in[0,1-\rho]$ and $\beta\notin\Bad(u,\delta)$.
If $S\subseteq D_{i+1}$ with $|S|\ge(1-\delta)n_{i+1}$ and $V\in\RS_{i+1}$ satisfy
$\Fold_\beta[w]=V$ on $S$, then there exists $G\in\RS_i$ with
\[
w=G\ \text{on}\ \fib(S):=\textstyle\bigcup_{y\in S}\fib(y)
\qquad\text{and}\qquad \Fold_\beta[G]=V .
\]
\end{lemma}

\begin{proof}
Since $\Fold_\beta[w]=\sum_j\beta^ju_j$ and $\beta$ is not bad, there are
$v_0,\dots,v_{a-1}\in\RS_{i+1}$ with $u_j[w]=v_j$ on $S$. Put
$G(X):=\sum_{j<a}X^jv_j(\sigma_{i+1}(X))$; then $\deg G<aL_{i+1}=L_i$ and $u_j[G]=v_j$ by
\cref{lem:components}(iii). For $y\in S$ the component vectors of $w$ and $G$ at $y$
coincide, hence $w=G$ on $\fib(y)$ by \cref{lem:components}(ii). Finally
$\Fold_\beta[G]=\sum_j\beta^jv_j\in\RS_{i+1}$ agrees with $V$ on $S$, and
$|S|\ge(1-\delta)n_{i+1}\ge\rho n_{i+1}=L_{i+1}$; two polynomials of degree $<L_{i+1}$
agreeing on $L_{i+1}$ points are equal.
\end{proof}

%=====================================================================
\section{Unique block decoding}
%=====================================================================

Put $\theta:=(1-\rho)/2$.

\begin{lemma}[Unique block decoding]\label{lem:unique-block}
Let $0\le i<k$ and $w:D_i\to\K$. There is at most one $G\in\K[X]_{<L_i}$ with
$\Db(w,G)<\theta$. If moreover $w$ takes values in $\F$, this $G$ lies in $\F[X]$.
\end{lemma}

\begin{proof}
If $\Db(w,G)<\theta$ and $\Db(w,G')<\theta$, then $G$ and $G'$ agree on more than
$(1-2\theta)n_{i+1}=\rho n_{i+1}=L_{i+1}$ fibres, contradicting
\cref{lem:codeword-fibres} unless $G=G'$.
For the second claim let $\phi$ be the $q$-power Frobenius of $\K/\F$, acting on
coefficients. Since $D_i\subseteq\F$ and $w$ is $\F$-valued,
$\phi(G)(x)=\phi(G(x))$ and $\phi(w(x))=w(x)$ for $x\in D_i$, so
$w(x)=G(x)\iff w(x)=\phi(G)(x)$. Hence $\Db(w,\phi(G))=\Db(w,G)<\theta$ and
$\phi(G)=G$ by uniqueness, i.e.\ $G\in\F[X]$.
\end{proof}

%=====================================================================
\section{The unique-decoding regime}
%=====================================================================

The only imported statement of this section is the proximity gap for parameterized
curves of Ben-Sasson, Carmon, Ishai, Kopparty and Saraf.

\begin{theorem}[Proximity gaps for curves, unique decoding {\cite[Thm.~1.5]{BCIKS20}}]
\label{thm:bciks-ud}
Let $a\ge2$, $u_0,\dots,u_{a-1}:E\to\K$ and $\delta\in[0,(1-\rho)/2)$. If
\[
\#\Bigl\{\beta\in\K:\ \Delta\Bigl(\textstyle\sum_{j<a}\beta^ju_j,\ C\Bigr)\le\delta\Bigr\}>(a-1)\,n,
\]
then there exist $S\subseteq E$ with $|S|\ge(1-\delta)n$ and $v_0,\dots,v_{a-1}\in C$ with
$u_j|_S=v_j|_S$ for all $j$.
\end{theorem}

\begin{proposition}[MCA in the unique-decoding regime]\label{prop:mca-ud}
For $0\le i<k$, $a=a_{i+1}$ and every $\delta\in[0,\theta)$,
\[
\eps_{\MCA}(\RS_{i+1},a,\delta)\le\frac{(a-1)\,n_{i+1}}{|\K|}.
\]
More precisely, for every $w:D_i\to\K$ there is a set $\mathcal E(w)\subseteq\K$ with
$|\mathcal E(w)|\le(a-1)n_{i+1}$, depending only on $w$, such that
$\Bad(u[w],\delta)\subseteq\mathcal E(w)$ for \emph{all} $\delta<\theta$ simultaneously.
\end{proposition}

\begin{proof}
By \cref{lem:components}(ii), every $u\in(\K^{D_{i+1}})^a$ equals $u[w]$ for a unique
$w:D_i\to\K$, so it suffices to treat $u=u[w]$. Write $n:=n_{i+1}$ and
$C:=\RS_{i+1}$.

\emph{Case 1: $\Db(w,\RS_i)\ge\theta$.} Let $\mathcal E(w):=\{\beta:\Delta(\Fold_\beta[w],C)<\theta\}$.
Every bad $\beta$ for $(u,\delta)$ with $\delta<\theta$ lies in $\mathcal E(w)$, since
condition (i) of \cref{def:mca} gives $\Delta(\Fold_\beta[w],C)\le\delta$. Suppose
$|\mathcal E(w)|>(a-1)n$. The values $\Delta(\Fold_\beta[w],C)$ lie in
$\frac1n\Z$; let $\delta^\ast$ be the largest element of $\frac1n\Z$ below $\theta$. Then
$\mathcal E(w)=\{\beta:\Delta(\Fold_\beta[w],C)\le\delta^\ast\}$ and $\delta^\ast<\theta$,
so \cref{thm:bciks-ud} yields $S$ with $|S|\ge(1-\delta^\ast)n$ and $v_j\in C$ with
$u_j[w]=v_j$ on $S$. As in the proof of \cref{lem:mca-fold}, $G:=\sum_jX^jv_j(\sigma_{i+1}(X))$
satisfies $w=G$ on $\fib(S)$, so $\Db(w,G)\le\delta^\ast<\theta$, a contradiction. Hence
$|\mathcal E(w)|\le(a-1)n$.

\emph{Case 2: $\Db(w,\RS_i)<\theta$.} Let $G$ be the unique polynomial of
\cref{lem:unique-block}, and let $Y:=\{y\in D_{i+1}: w|_{\fib(y)}\ne G|_{\fib(y)}\}$, so
$|Y|<\theta n$. For $y\in Y$ the polynomial
\[
p_y(\beta):=\sum_{j<a}\beta^j\bigl(u_j[w](y)-u_j[G](y)\bigr)\in\K[\beta]
\]
is nonzero by \cref{lem:components}(ii) and has degree $\le a-1$. Let
$\mathcal E(w):=\bigcup_{y\in Y}p_y^{-1}(0)$, so $|\mathcal E(w)|\le(a-1)|Y|<(a-1)n$.

Let $\beta\notin\mathcal E(w)$ and let $S,v$ be as in \cref{def:mca}(i) with
$\delta<\theta$. Then $\Fold_\beta[w](y)\ne\Fold_\beta[G](y)$ exactly for $y\in Y$.
The codewords $v$ and $\Fold_\beta[G]$ (the latter in $C$ by \cref{lem:fold-poly})
can disagree only on $(D_{i+1}\setminus S)\cup Y$, a set of size
$<\delta n+\theta n<2\theta n=(1-\rho)n$. So they agree on more than $\rho n=L_{i+1}$
points and $v=\Fold_\beta[G]$. On $S$ we then have
$\Fold_\beta[w]=v=\Fold_\beta[G]$, hence $S\cap Y=\emptyset$, hence
$u_j[w]=u_j[G]$ on $S$ for all $j$, and $(u_j[G])_j\in C^a$ by \cref{lem:components}(iii).
Condition (ii) of \cref{def:mca} fails, so $\beta$ is not bad.
\end{proof}

\begin{corollary}[Accumulated unique-decoding error]\label{cor:mca-ud-sum}
For every $\delta<\theta$,
\[
\sum_{i=0}^{k-1}\eps_{\MCA}(\RS_{i+1},a_{i+1},\delta)\le\frac{N-B}{|\K|}.
\]
\end{corollary}

\begin{proof}
$(a_{i+1}-1)n_{i+1}=n_i-n_{i+1}$ by \cref{lem:fibres}, and
$\sum_{i<k}(n_i-n_{i+1})=n_0-n_k=N-B$.
\end{proof}

%=====================================================================
\section{The Johnson regime}
%=====================================================================

The univariate powers generator is a polynomial generator in the sense of Bordage,
Chiesa, Guan and Manzur, who prove mutual correlated agreement for all polynomial
generators and Reed--Solomon codes up to the Johnson radius~\cite{BCGM26}. We record
the statement in the form used here.

\begin{theorem}[MCA up to the Johnson radius {\cite{BCGM26}}]\label{thm:mca-johnson}
For every $a\ge2$ and $\eta\in(0,1-\sqrt\rho)$ there is an explicit quantity
$\eps^{J}_{\MCA}(n,\rho,a,\eta,|\K|)$ such that
\[
\eps_{\MCA}\bigl(\RS[\K,E,\ell],a,\,1-\sqrt\rho-\eta\bigr)\le\eps^{J}_{\MCA}(n,\rho,a,\eta,|\K|).
\]
\end{theorem}

\begin{remark}[Instantiation of $\eps^J_{\MCA}$]\label{rem:johnson-constants}
The parameter sets of the parameter chapter use $\eps^J_{\MCA}$ only through the sum
$\sum_{i<k}\eps^J_{\MCA}(n_{i+1},\rho,a_{i+1},\eta,|\K|)$. The original correlated
agreement bound of~\cite{BCIKS20} in this regime is quadratic in $n$; the bounds
of~\cite{BCHKS25} (correlated agreement) and~\cite{DKT26} (mutual correlated agreement)
are linear in $n$ with polynomial dependence on $1/\eta$, and these are the relevant
ones for $|\K|\approx2^{192}$ and $n\le2^{24}$. The explicit constants are inserted in the parameter chapter.
\end{remark}

%=====================================================================
\section{Block lists and the Johnson list bound}
%=====================================================================

\begin{definition}[Block list]\label{def:block-list}
For $w:D_0\to\K$ and $\alpha\in(0,1]$,
\[
\Lambda_\alpha(w):=\bigl\{G\in\K[X]_{<L}:\ \agr(w,G)\ge\alpha\bigr\},\qquad
\ell_\alpha:=\max_{w}|\Lambda_\alpha(w)| .
\]
\end{definition}

\begin{proposition}[List sizes]\label{prop:list-size}
\begin{enumerate}[label=\textup{(\roman*)}]
\item If $\alpha>\frac{1+\rho}2$ then $\ell_\alpha\le1$.
\item If $\alpha>\sqrt\rho$ then
\[
\ell_\alpha\le\frac{\alpha-\rho}{\alpha^2-\rho}.
\]
\end{enumerate}
\end{proposition}

\begin{proof}
(i) is \cref{lem:unique-block} with $i=0$, since $\agr>\frac{1+\rho}2$ means
$\Db<\theta$.

(ii) Write $n:=n_1$ and let $G_1,\dots,G_M\in\Lambda_\alpha(w)$ be distinct. Let
$t:=\lceil\alpha n\rceil$ and choose sets $S_r\subseteq D_1$ with $|S_r|=t$ such that
$w=G_r$ on $\fib(y)$ for every $y\in S_r$. For $y\in D_1$ let
$c_y:=\#\{r: y\in S_r\}$. Then $\sum_yc_y=Mt$. If $y\in S_r\cap S_{r'}$ with $r\ne r'$,
then $G_r=w=G_{r'}$ on $\fib(y)$; by \cref{lem:codeword-fibres} each pair $r\ne r'$ has
fewer than $\rho n$ such $y$, so
\[
\sum_yc_y(c_y-1)=\sum_{r\ne r'}|S_r\cap S_{r'}|\le M(M-1)\rho n .
\]
By Cauchy--Schwarz $\sum_yc_y^2\ge(Mt)^2/n$. Hence
$(Mt)^2/n-Mt\le M(M-1)\rho n$; dividing by $Mn$ and writing $x:=t/n\ge\alpha$,
\[
M(x^2-\rho)\le x-\rho .
\]
Since $x\ge\alpha>\sqrt\rho$, $x^2-\rho>0$ and $M\le\frac{x-\rho}{x^2-\rho}$. The function
$x\mapsto\frac{x-\rho}{x^2-\rho}$ is decreasing on $(\sqrt\rho,1]$: its derivative has
numerator $-(x^2-2\rho x+\rho)=-\bigl((x-\rho)^2+\rho(1-\rho)\bigr)<0$. Therefore
$M\le\frac{\alpha-\rho}{\alpha^2-\rho}$.
\end{proof}

\begin{remark}\label{rem:list-galois}
If $w$ is $\F$-valued, the argument of \cref{lem:unique-block} shows that
$\Lambda_\alpha(w)$ is stable under the Frobenius of $\K/\F$. In the regime of
\cref{prop:list-size}(i) its unique element is therefore in $\F[X]$; in the regime of
\cref{prop:list-size}(ii) elements of the list may have coefficients in $\K\setminus\F$.
\end{remark}
